% Canonical synthesis and extension of the reflected Stieltjes calculus.
\section{Reflection closes the Stieltjes convolution algebra}
\label{stref:section}
The source report \cite{stconv:calculus-report} already identifies the joint algebra of both orientations. We restore that structural result here before deriving a pole-free many-factor generator and its neutral-order identities. The one-generator algebra in Theorem~\ref{stconv:thm:Bell} has a
finite extension when reflection is allowed. This extension gives one
normal form for every finite convolution of generalized Stieltjes
functions with independently chosen orientations. Its derivation uses
the classical Hurwitz Fourier formula \cite[Eq.~25.11.9]{intdist:DLMFHurwitz};
no literature-priority or independence claim for numerical constants
is intended.

Let $I=\delta_0-1$, let $\check T(x)=T(-x)$ denote periodic reflection,
and put $X=G_0$, $Y=\check G_0$. The unit in every algebra below is $I$.
For a polynomial in $X,Y$, multiplication means convolution.

\begin{theorem}[The exact reflection algebra]\label{stref:algebra}
The complex convolution algebra generated by all $G_n$ and
$\check G_n$, $n\ge0$, is
\[
 \mathcal A\cong
 \mathbb C[x,y]/\big((x-y)^2+\pi^2\big)
 \cong\mathbb C[t]\times\mathbb C[t].
\]
The displayed quadratic relation is the entire ideal of relations.
Equivalently, with $D=X-Y$, every element has a unique form
$P(X)+D*Q(X)$, with $P,Q\in\mathbb C[t]$.
A second basis, adapted to generalized Stieltjes functions, is
\[
 I,\quad G_0,\check G_0,\quad G_1,\check G_1,\quad G_2,\check G_2,\ \ldots .
\]
The two spectral idempotents are
\[
 E_+=\frac12\left(I-\frac{D}{i\pi}\right),\qquad
 E_-=\frac12\left(I+\frac{D}{i\pi}\right).
\]
They satisfy $E_\pm*E_\pm=E_\pm$, $E_+*E_-=0$,
$E_++E_-=I$. Their Fourier coefficients are respectively the indicators
of strictly positive and strictly negative integers. Reflection and
complex conjugation each exchange them.
\end{theorem}
\begin{proof}
At $k\ne0$, the established Fourier coefficients are
\[
 \widehat X(k)=-\gamma-\log(2\pi|k|)-\frac{i\pi}{2}\operatorname{sgn}k,
 \qquad
 \widehat Y(k)=-\gamma-\log(2\pi|k|)+\frac{i\pi}{2}\operatorname{sgn}k.
\]
Thus $\widehat D(k)=-i\pi\operatorname{sgn}k$ and
$D*D=-\pi^2I$, including the zero coefficient. The Bell formula
expresses every $G_n$ as a polynomial in $X$, and its reflection
expresses every $\check G_n$ as the same polynomial in $Y$.
Polynomial division by $D^2+\pi^2$ gives the first normal form.

If $P(X)+D*Q(X)=0$, positive Fourier frequencies give
$P(t)-i\pi Q(t)=0$ at infinitely many distinct values
$t=-\gamma-\log(2\pi k)-i\pi/2$. Hence that polynomial is zero.
Negative frequencies likewise imply $P+i\pi Q=0$ identically.
Therefore $P=Q=0$; there are no further relations. Evaluation at
$D=-i\pi$ and $D=i\pi$ proves the product-of-polynomial-algebras
assertion and the idempotent statements.

For the Stieltjes basis, the Bell polynomial has leading term
$G_n=X^{*(n+1)}/(n+1)$ and no term of degree $n$.
Modulo $D^2=-\pi^2I$, the difference
$G_n-\check G_n$ has leading term $D*X^{*n}$.
Thus triangular elimination in the basis
$I,X,D,X^{*2},D*X,\ldots$ proves both spanning and independence.
The same proof over $\mathbb R$ gives
$\mathbb R[x,D]/(D^2+\pi^2)$, which is isomorphic to
$\mathbb C[t]$ as a real algebra; the spectral idempotents themselves
belong to the complexification.
\end{proof}

\begin{theorem}[All orientations in one analytic generator]\label{stref:generator}
Set
\[
 F(u)=I+u\sum_{n\ge0}G_n\frac{u^n}{n!},\qquad
 \check F(u)=I+u\sum_{n\ge0}\check G_n\frac{u^n}{n!}.
\]
Choose any $p,q\ge0$ with $p+q\ge1$, and put
\[
 A=\sum_{i=1}^p u_i,\qquad B=\sum_{j=1}^q v_j,\qquad S=A+B,
 \qquad Q=\frac{\prod_i\Gamma(1+u_i)\prod_j\Gamma(1+v_j)}{\Gamma(1+S)}.
\]
As a holomorphic germ of periodic distributions at the origin,
\begin{align}
 \mathop{*}_{i=1}^pF(u_i)*\mathop{*}_{j=1}^q\check F(v_j)
 =Q\bigg\{&\frac{\cos(\pi(A-B)/2)}{\cos(\pi S/2)}I\notag\\
 &+\frac{S}{\sin\pi S}\left[
   \sin(\pi A)\sum_{n\ge0}G_n\frac{S^n}{n!}
  +\sin(\pi B)\sum_{n\ge0}\check G_n\frac{S^n}{n!}
                  \right]\bigg\}.
 \label{stref:master}
\end{align}
Here $S/\sin\pi S$ has value $1/\pi$ at zero. In particular,
there is no diagonal singularity at $A+B=0$.
\end{theorem}
\begin{proof}
The previous master identity is $F(u)=\Gamma(1+u)R_u$.
At a nonzero Fourier frequency the product on the left is
\[
 \prod_i\Gamma(1+u_i)\prod_j\Gamma(1+v_j)
 (2\pi|k|)^{-S}
 \exp\{-i\pi(A-B)\operatorname{sgn}(k)/2\}.
\]
When $\sin\pi S\ne0$, the elementary sine addition identity gives
\[
 R_A*\check R_B
  =\frac{\sin\pi A}{\sin\pi S}R_S
   +\frac{\sin\pi B}{\sin\pi S}\check R_S.
\]
Insert $R_S=F(S)/\Gamma(1+S)$ and use
\[
 \frac{\sin\pi A+\sin\pi B}{\sin\pi S}
  =\frac{\cos(\pi(A-B)/2)}{\cos(\pi S/2)}.
\]
This yields \eqref{stref:master} on an open set. The displayed right
side is holomorphic near the full origin, including $S=0$, so the
identity extends there. The Gamma factors and the Fourier coefficients
have locally uniform polynomial growth, and pairing with a smooth test
function therefore justifies every parameter expansion. All zero
Fourier coefficients vanish, as required by the unit $I$.
\end{proof}

\begin{corollary}[Finite all-index reduction and its leading terms]
\label{stref:allindices}
For fixed nonnegative indices $m_1,\ldots,m_p,n_1,\ldots,n_q$, define
\[
 d_+=\sum_i(m_i+1),\quad d_-=\sum_j(n_j+1),\quad d=d_++d_-,
 \qquad M=\prod_i(m_i+1)\prod_j(n_j+1).
\]
Then
\begin{align}
 &G_{m_1}*\cdots*G_{m_p}*
  \check G_{n_1}*\cdots*\check G_{n_q}\notag\\
 &\quad=\frac{d_+}{M}G_{d-1}+\frac{d_-}{M}\check G_{d-1}
       +cI+\sum_{j=0}^{d-3}(a_jG_j+b_j\check G_j).
 \label{stref:normal}
\end{align}
The sum is empty if $d<3$. For $d\ge2$, the coefficient of either orientation of
$G_{d-2}$ is zero. All coefficients belong to
$\mathbb Q[\zeta(2),\ldots,\zeta(d)]$.
They are uniquely determined by multiplying the coefficient of
$\prod_i u_i^{m_i+1}\prod_j v_j^{n_j+1}$ in
\eqref{stref:master} by $\prod_i m_i!\prod_j n_j!$.
Assigning weight $j+1$ to $G_j$ and $\check G_j$ and weight $k$ to
$\zeta(k)$ makes every term have weight $d$.
\end{corollary}
\begin{proof}
The coefficient extraction follows directly from the factors on the
left. Only total degree at most $d$ on the right can contribute, so
it is a finite algorithm. Explicitly,
\[
 \log Q=\sum_{r\ge2}\frac{(-1)^r\zeta(r)}r
  \left(\sum_i u_i^r+\sum_jv_j^r-S^r\right).
\]
The scalar trigonometric factor has only even-degree corrections to
its constant term, and
$(S/\sin\pi S)\sin\pi A=A+O((|A|+|B|)^3)$, with the analogous
formula for $B$. Hence the top terms arise from
$A S^{d-1}G_{d-1}/(d-1)!$ and its reflected counterpart, giving
$d_+/M$ and $d_-/M$. There is no degree-one correction, so
$G_{d-2}$ cannot occur. All trigonometric coefficients are rational
polynomials in $\pi^2=6\zeta(2)$. The Gamma expansion proves the
remaining coefficient and weight assertions. Uniqueness follows from
Theorem~\ref{stref:algebra}.
\end{proof}

For example, the first asymmetric triple product is
\begin{equation}\label{stref:triple}
 X^{*2}*Y=2G_2+\check G_2+2\zeta(2)G_0+\zeta(2)\check G_0
                 +2\zeta(3)I.
\end{equation}
Reflection gives $X*Y^{*2}$. These identities retain the contact mass;
replacing $I$ by $\delta_0$ changes their smooth values as well as
their means.

Away from the integer singularity, replace $G_j$ and $\check G_j$
by $\gamma_j(x)$ and $\gamma_j(1-x)$, and replace $I$ by $-1$.
An ordinary-integral realization can be built recursively using the
absolutely convergent subtractions in Theorem~\ref{stconv:thm:integral}
and the reflected correlation formula \eqref{shj:Jregular}.
For instance $G_1*\check G_0=J_{01}$ and
\eqref{stref:triple} is equivalently
\[
 2J_{01}(x)-\zeta(2)\gamma_0(1-x)
  =2\gamma_2(x)+\gamma_2(1-x)+2\zeta(2)\gamma_0(x)
    +\zeta(2)\gamma_0(1-x)-2\zeta(3),\quad 0<x<1.
\]
This last scalar identity is also an explicit overlap check against the
previous pair theorem; the extension here is the single all-orientation
generator and its coefficient laws. Arbitrary pointwise products with
coincident singularities are not asserted to exist, and no unregularized
multiple integral is implicit.

\subsection{An alternating correlation identity at every index}
The diagonal $A+B=0$ in the pole-free generator is especially simple.
For $F$ as above, Euler reflection gives the exact germ
\begin{equation}\label{stref:opposite-orders}
 F(u)*\check F(-u)=\pi u\cot(\pi u)I+u(X-Y).
\end{equation}
Indeed \eqref{stref:master} with $A=u$, $B=-u$ has bracket
$\cos(\pi u)I+\sin(\pi u)(X-Y)/\pi$, while
$\Gamma(1+u)\Gamma(1-u)=\pi u/\sin\pi u$.
Comparing coefficients proves, for every $k\ge0$,
\begin{equation}\label{stref:alternating}
 \sum_{j=0}^k(-1)^{k-j}\binom{k}{j}G_j*\check G_{k-j}
 =\frac{G_{k+1}+(-1)^k\check G_{k+1}}{k+1}
   +\mathbf1_{2\mid k}\,2k!\zeta(k+2)I.
\end{equation}
Thus all lower Stieltjes layers cancel in this alternating sum.
The equivalent scalar identity for the fixed normalization
\eqref{shj:Jregular}, involving only ordinary subtracted integrals, is
\begin{equation}\label{stref:alternating-scalar}
 \boxed{\displaystyle
 \sum_{j=0}^k(-1)^j\binom{k}{j}J_{j,k-j}(a)
 =\frac{\gamma_{k+1}(a)+(-1)^k\gamma_{k+1}(1-a)}{k+1}
  -\mathbf1_{2\mid k}\,2k!\zeta(k+2),\quad 0<a<1.}
\end{equation}
In particular,
\[
 J_{02}-2J_{11}+J_{20}
 =\frac{\gamma_3(a)+\gamma_3(1-a)}3-4\zeta(4),
\]
whereas
\[
 J_{03}-3J_{12}+3J_{21}-J_{30}
 =\frac{\gamma_4(a)-\gamma_4(1-a)}4.
\]
\eqref{stref:opposite-orders} is an analytic germ by
Theorem~\ref{stref:generator}, so coefficient extraction followed by
the established separated-singularity restriction proves the scalar
formula without requiring an ordinary integral before regularization.

\begin{corollary}[Many-factor cancellation of the spectral order]
\label{stref:neutral-many}
On the hyperplane $S=\sum_i u_i+\sum_jv_j=0$, the exact germ is
\begin{equation}\label{stref:neutral-many-eq}
 \mathop{*}_iF(u_i)*\mathop{*}_j\check F(v_j)
 =\left(\prod_i\Gamma(1+u_i)\prod_j\Gamma(1+v_j)\right)
   \left(\cos(\pi A)I+\frac{\sin(\pi A)}\pi D\right).
\end{equation}
Thus every Taylor coefficient after any linear parametrization of this
hyperplane lies in the two-dimensional span of $I$ and $D$; all powers
of the logarithmic Fourier variable cancel. This holds for any number
of factors. It concerns convolutions of the normalized spectral
generators $F$, rather than arbitrary pointwise products of Stieltjes
functions.
\end{corollary}
\begin{proof}
Set $S=0$ in the pole-free germ \eqref{stref:master}; the only surviving
Stieltjes coefficient is $G_0$ or $\check G_0$, and $B=-A$.
Equivalently the magnitude factor $(2\pi|k|)^{-S}$ in every nonzero
Fourier coefficient is one. This also proves the assertion on Taylor
coefficients directly, with no limit along an unspecified path.
\end{proof}
